\section{Arquitectura del Sistema y Diagrama de Despliegue UML}
\label{sec:arquitectura_sistema}
% =============================================================================

\subsection{Arquitectura en Capas del Sistema y Ubicación de Modelos UWE}

Para garantizar una alta cohesión, bajo acoplamiento y facilidad de mantenimiento, el prototipo se estructura según un modelo de arquitectura multicapa de seis niveles, como se ilustra en la Figura \ref{fig:uml_arquitectura}. Los artefactos de ingeniería web UWE gobiernan directamente las dos capas superiores (Presentación/Navegación y Control de Procesos REST), acoplándose con las capas de dominio electoral, abstracción FFI, núcleo nativo en C y persistencia relacional.

\begin{figure}[!htbp]
	\centering
	\resizebox{\textwidth}{!}{%
	\begin{tikzpicture}[
		layer/.style={rectangle, draw=black, fill=white, line width=0.8pt, rounded corners=1.5pt, minimum width=13.2cm, minimum height=1.1cm, font=\scriptsize, align=center},
		subcomp/.style={rectangle, draw=black, fill=white, line width=0.5pt, rounded corners=1pt, font=\tiny, align=center, inner sep=3pt},
		dep/.style={-{Latex[length=2mm, width=1.4mm]}, line width=0.8pt, draw=black}
	]
		% Capa 1: Presentación
		\node[layer] (capa1) at (0, 0) {
			\textbf{Capa 1: Presentación y Navegación Web UWE (Cliente de Votación)}\\
		};
		\node[subcomp] at (-4.3, -0.18) {Cabina Web UWE (\texttt{«presentation page»})};
		\node[subcomp] at (0.0, -0.18) {Panel Registro AR (\texttt{«input form»})};
		\node[subcomp] at (4.3, -0.18) {Portal Ciudadano Auditoría (\texttt{«query»})};

		% Capa 2: Servicios Web y API REST
		\node[layer, below=0.55cm of capa1] (capa2) {
			\textbf{Capa 2: Servicios Web, Controlador de Procesos UWE y API REST (Flask WSGI)}\\
		};
		\node[subcomp] at (-4.5, -1.83) {Endpoints: \texttt{/api/votar}, \texttt{/api/escrutinio}};
		\node[subcomp] at (0.2, -1.83) {Serializador Canónico (\texttt{canonical\_json})};
		\node[subcomp] at (4.6, -1.83) {Guarda de Estado OCL y Sesión};

		% Capa 3: Lógica del Dominio Electoral
		\node[layer, below=0.55cm of capa2] (capa3) {
			\textbf{Capa 3: Lógica del Dominio y Modelo de Contenido UWE (Python)}\\
		};
		\node[subcomp] at (-4.5, -3.48) {\texttt{VotantePQC}: Formulación y Cifrado};
		\node[subcomp] at (-1.2, -3.48) {\texttt{AutoridadRegistro}: Padrón};
		\node[subcomp] at (2.0, -3.48) {\texttt{ServidorVotacion}: Urna};
		\node[subcomp] at (4.8, -3.48) {\texttt{AutoridadConteo}: Tally};

		% Capa 4: Abstracción Criptográfica FFI
		\node[layer, below=0.55cm of capa3] (capa4) {
			\textbf{Capa 4: Abstracción Criptográfica y Enlace FFI (ctypes)}\\
		};
		\node[subcomp] at (-4.2, -5.13) {\texttt{pqc\_wrapper.py}: Wrappers ML-KEM/DSA};
		\node[subcomp] at (0.2, -5.13) {\texttt{crypto\_utils.py}: AES-256-GCM / SHA3-256};
		\node[subcomp] at (4.4, -5.13) {Manejador de Punteros \texttt{c\_uint8}};

		% Capa 5: Núcleo Nativo C Core
		\node[layer, below=0.55cm of capa4] (capa5) {
			\textbf{Capa 5: Núcleo Criptográfico Nativo en Lenguaje C (PQClean / libpqc\_core.dylib)}\\
		};
		\node[subcomp] at (-3.8, -6.78) {ML-KEM-768 (FIPS 203): KeyGen, Encaps, Decaps};
		\node[subcomp] at (3.5, -6.78) {ML-DSA-65 (FIPS 204): KeyGen, Sign, Verify};

		% Capa 6: Persistencia
		\node[layer, below=0.55cm of capa5] (capa6) {
			\textbf{Capa 6: Persistencia y Almacenamiento Relacional Disociado (PostgreSQL / SQLite)}\\
		};
		\node[subcomp] at (-4.2, -8.43) {Padrón Electoral (\texttt{tbl\_padron})};
		\node[subcomp] at (0.0, -8.43) {Ánfora Digital Cifrada (\texttt{tbl\_anfora})};
		\node[subcomp] at (4.2, -8.43) {Bitácora y Recibos (\texttt{tbl\_recibos})};

		% Conectores limpios entre capas consecutivas
		\draw[dep] (capa1) -- (capa2);
		\draw[dep] (capa2) -- (capa3);
		\draw[dep] (capa3) -- (capa4);
		\draw[dep] (capa4) -- (capa5);
		\draw[dep] (capa3.east) to[bend left=45] (capa6.east);
	\end{tikzpicture}
	}%
	\caption{Diagrama de Arquitectura en Capas del sistema de votación post-cuántica articulado con UWE.}
	\label{fig:uml_arquitectura}
	\par\vspace{3pt}{\footnotesize\textit{Fuente: Elaboración propia.}}
\end{figure}

Cada capa mantiene responsabilidades rigurosamente delimitadas:
\begin{itemize}
	\item \textbf{Capa 1 (Presentación y Navegación Web UWE):} Materializa los modelos de Navegación y Presentación de UWE. Provee las interfaces visuales responsivas para la cabina del votante, el panel de registro de la AR y el portal ciudadano de auditoría, aplicando en el cliente las transiciones entre vistas (\texttt{«view port»}) y las guardas de navegación.
	\item \textbf{Capa 2 (Servicios Web, Procesos UWE y API REST):} Expone los puntos de entrada REST y orquesta las clases de proceso UWE (\texttt{«process class»}). Verifica las precondiciones OCL de sesión e implementa la serialización determinista canónica antes del procesamiento criptográfico.
	\item \textbf{Capa 3 (Lógica del Dominio y Contenido UWE):} Instancia las clases del Modelo de Contenido UWE y coordina a los cuatro actores electorales: validación de credenciales, verificación de firmas digitales, detección de intentos de doble votación y escrutinio final.
	\item \textbf{Capa 4 (Abstracción FFI Ctypes):} Realiza la conversión segura de tipos entre las estructuras de Python y los buffers contiguos de memoria en lenguaje C (\texttt{c\_uint8}).
	\item \textbf{Capa 5 (Núcleo C Core):} Ejecuta el procesamiento sobre retículos algebraicos (multiplicación polinomial en el dominio NTT, convolución negacíclica y expansión pseudoaleatoria Keccak SHAKE-128/256), compilado con optimizaciones nativas para la arquitectura destino.
	\item \textbf{Capa 6 (Persistencia Relacional):} Almacena de forma transaccional y normalizada en 3FN los registros electorales, manteniendo la disociación estructural entre el padrón y el ánfora de boletas cifradas.
\end{itemize}

\subsection{Diagrama de Despliegue UML del Entorno de Ejecución}

La Figura \ref{fig:uml_despliegue} presenta el Diagrama de Despliegue UML del sistema, modelando los nodos de hardware, los entornos de ejecución, los artefactos de software instalados y los canales de comunicación con la Autoridad de Conteo aislada físicamente fuera de línea (\textit{air-gapped}).

\begin{figure}[!htbp]
	\centering
	\resizebox{\textwidth}{!}{%
	\begin{tikzpicture}[
		node distance=1.0cm and 1.2cm,
		nodebox/.style={rectangle, draw=black, fill=white, line width=0.8pt, rounded corners=1.5pt, minimum width=3.8cm, minimum height=3.0cm, font=\scriptsize, align=center},
		artbox/.style={rectangle, draw=black, fill=white, line width=0.5pt, rounded corners=1pt, font=\tiny, align=center, inner sep=3.5pt},
		netlink/.style={line width=0.8pt, draw=black, font=\tiny}
	]
		% Nodo 1: Terminal Cliente
		\node[nodebox] (client) at (-5.2, 3.2) {
			\guillemotleft device\guillemotright\\
			\textbf{Terminal Votante / Cabina}\\
			\vspace{1.8cm}
		};
		\node[artbox, below=0.4cm of client.north] {\guillemotleft execution env\guillemotright\\Navegador Web / Interfaz UWE};
		\node[artbox, below=1.3cm of client.north] {\guillemotleft artifact\guillemotright\\PQC Client Engine (Wasm / Python)};

		% Nodo 2: Servidor Web SV
		\node[nodebox, minimum width=4.6cm] (server) at (1.2, 3.2) {
			\guillemotleft device\guillemotright\\
			\textbf{Servidor Electoral Central (SV)}\\
			\vspace{1.8cm}
		};
		\node[artbox, below=0.4cm of server.north] {\guillemotleft execution env\guillemotright\\Reverse Proxy Nginx + Gunicorn};
		\node[artbox, below=1.3cm of server.north] {\guillemotleft artifact\guillemotright\\Flask Web API + \texttt{libpqc\_core.dylib}\\(Clave pública $ek_{\text{AC}}$ publicada)};

		% Nodo 3: Servidor de Base de Datos
		\node[nodebox] (db) at (6.8, 3.2) {
			\guillemotleft device\guillemotright\\
			\textbf{Servidor Base Datos}\\
			\vspace{1.8cm}
		};
		\node[artbox, below=0.4cm of db.north] {\guillemotleft DBMS\guillemotright\\PostgreSQL / SQLite3};
		\node[artbox, below=1.3cm of db.north] {\guillemotleft artifact\guillemotright\\Esquema Relacional 3FN\\(Volumen Cifrado LUKS)};

		% Nodo 4: Estación AR
		\node[nodebox] (ar_node) at (-5.2, -2.0) {
			\guillemotleft device\guillemotright\\
			\textbf{Estación de Registro (AR)}\\
			\vspace{1.8cm}
		};
		\node[artbox, below=0.4cm of ar_node.north] {\guillemotleft execution env\guillemotright\\Terminal de Empadronamiento};
		\node[artbox, below=1.3cm of ar_node.north] {\guillemotleft artifact\guillemotright\\Módulo AR + Clave $sk_{\text{AR}}$};

		% Nodo 5: Estación AC (Air Gapped)
		\node[nodebox, minimum width=4.6cm] (ac_node) at (1.2, -2.0) {
			\guillemotleft device\guillemotright\\
			\textbf{Estación Conteo (AC - Air-Gapped)}\\
			\vspace{1.8cm}
		};
		\node[artbox, below=0.4cm of ac_node.north] {\guillemotleft secure element\guillemotright\\Token Criptográfico / HSM Seguro};
		\node[artbox, below=1.3cm of ac_node.north] {\guillemotleft artifact\guillemotright\\Módulo Escrutinio + Clave $dk_{\text{AC}}$\\(Genera $dk_{\text{AC}}$, exporta $ek_{\text{AC}}$)};

		% Conexiones de red
		\draw[netlink] (client) -- (server) 
		node[midway, above] {HTTPS / TLS 1.3}
		node[midway, below] {(Puerto 443 / 5000)};
		
		\draw[netlink] (server) -- (db) 
		node[midway, above] {TCP/IP 5432}
		node[midway, below] {(Red Interna Aislada)};
		
		\draw[netlink, dashed] (ar_node) -- (client)
		node[midway, left] {Canal Seguro / Presencial}
		node[midway, right] {(Emisión Credencial)};
		
		\draw[netlink, dashed] (ac_node) -- (server)
		node[midway, left] {1. Exportación de $ek_{\text{AC}}$};
		\draw[netlink, dashed] (server) -- (ac_node)
		node[midway, right] {2. Entrega de Lote Permutado (Post-Cierre)};
	\end{tikzpicture}
	}%
	\caption{Diagrama de Despliegue UML del entorno de hardware, comunicaciones y custodia física de $dk_{\text{AC}}$.}
	\label{fig:uml_despliegue}
	\par\vspace{3pt}{\footnotesize\textit{Fuente: Elaboración propia.}}
\end{figure}

El esquema de despliegue contempla tres garantías operativas fundamentales:
\begin{enumerate}
	\item \textbf{Canal de Votación Abierto bajo Modelo Dolev--Yao:} Los clientes web se comunican con el Servidor de Votación mediante HTTPS (TLS 1.3). Siguiendo las directrices del modelo Dolev--Yao, la seguridad de la boleta no se confía a la capa de transporte: la confidencialidad e integridad están respaldadas criptográficamente de extremo a extremo por ML-KEM-768 y ML-DSA-65.
	\item \textbf{Aislamiento Físico de la Autoridad de Conteo (\textit{Air-Gapped}):} La estación de escrutinio opera desconectada físicamente de cualquier red. La clave privada de desencapsulamiento $dk_{\text{AC}}$ se genera en su interior y nunca abandona el entorno protegido; únicamente la clave pública $ek_{\text{AC}}$ se exporta para su publicación, y tras el cierre comicial, el lote permutado de boletas anónimas se transfiere fuera de banda para el conteo.
	\item \textbf{Persistencia Segura:} El motor relacional se aloja en un segmento de red privado no enrutable hacia el exterior, almacenando los datos en volúmenes protegidos mediante cifrado de disco a nivel de bloque.
\end{enumerate}

\subsection{Diagrama de Componentes de Software UML}
\label{sec:diagrama_componentes}

La Figura \ref{fig:uml_componentes} expone el Diagrama UML de Componentes de Software, modelando los paquetes, las dependencias de compilación y las interfaces de enlace nativo entre la capa web UWE y el núcleo en C.

\begin{figure}[!htbp]
	\centering
	\begin{tikzpicture}[
		node distance=0.8cm and 1.2cm,
		comp/.style={rectangle, draw=black, fill=white, line width=0.8pt, rounded corners=1.5pt, minimum width=4.0cm, minimum height=0.9cm, font=\scriptsize, align=center},
		dep/.style={-{Latex[length=2mm, width=1.4mm]}, line width=0.7pt, dashed, draw=black}
	]
		% Capa 1: C Nativo
		\node[comp] (c_core) at (0, 0) {
			\guillemotleft component\guillemotright\\
			\textbf{PQClean C Core}\\
			\textit{ML-KEM-768 / ML-DSA-65}
		};
		\node[comp] (c_dyn) at (6.5, 0) {
			\guillemotleft artifact\guillemotright\\
			\textbf{libpqc\_core.dylib}\\
			\textit{Binario compilado ARM64/x86}
		};
		
		% Capa 2: FFI Wrapper & Utils
		\node[comp] (wrapper) at (0, -2.1) {
			\guillemotleft component\guillemotright\\
			\textbf{Python FFI Wrapper}\\
			\textit{pqc\_wrapper.py (ctypes)}
		};
		\node[comp] (utils) at (6.5, -2.1) {
			\guillemotleft component\guillemotright\\
			\textbf{Crypto Utilities}\\
			\textit{AES-GCM, SHA3, canonical\_json}
		};
		
		% Capa 3: Actores
		\node[comp] (actores) at (3.25, -4.2) {
			\guillemotleft component\guillemotright\\
			\textbf{Módulos de Entidades Electorales}\\
			\textit{Votante, AR, Servidor, Conteo (Tally)}
		};
		
		% Capa 4: Orquestador & UI
		\node[comp] (orch) at (0, -6.3) {
			\guillemotleft component\guillemotright\\
			\textbf{Orquestador E2E}\\
			\textit{protocolo\_pqc.py}
		};
		\node[comp] (web) at (6.5, -6.3) {
			\guillemotleft component\guillemotright\\
			\textbf{Capa Web UWE \& REST API}\\
			\textit{Flask app.py + Vistas UWE}
		};
		
		% Conexiones
		\draw[dep] (c_core) -- (c_dyn) node[midway, above, font=\tiny\itshape] {compilación gcc/clang};
		\draw[dep] (wrapper) -- (c_dyn) node[midway, right, font=\tiny\itshape] {dlopen / ctypes};
		\draw[dep] (actores) -- (wrapper);
		\draw[dep] (actores) -- (utils);
		\draw[dep] (orch) -- (actores);
		\draw[dep] (web) -- (actores);
	\end{tikzpicture}
	\caption{Diagrama UML de Componentes de Software de la implementación PQC.}
	\label{fig:uml_componentes}
	\par\vspace{3pt}{\footnotesize\textit{Fuente: Elaboración propia.}}
\end{figure}
